% 05_discussion.tex -- discussion and conclusion, draft 23 Oct 2026.
\section{Discussion}
\label{sec:discussion}

\paragraph{Rigid links in the twin, slack tethers on the rig}
The twin models every cable as a rigid link, which can push as well as pull and welds wherever its end lies. This
flatters the rejoin and also creates its own failure: on a T, the rejoiner's rigid link set the ring's pitch, so a
rising height reference lifted the ring's free side, and the height hold that cured it is a patch for the twin's
links. On the rig the tethers go slack, so that lift cannot arise, but the attach is harder: the magnet catches only on
contact, which is why the rig's attach is confirmed from motion capture before the newcomer is folded in. The rig's
release reaction, by contrast, needs no such caveat: it was measured on the rig, and the twin's release on longer
cables agreed in resize time and tilt.

\paragraph{Timing}
Every reconfiguration in Section~\ref{sec:results} is paced by detection and blending rather than by the solver: a
lost carrier is seen from its cable's stretch in 0.33--0.42\,s and the problem is resized within about 0.1\,s more,
and a rejoin's fold-in is blended over 10\,s. The solver sets the size of the team instead. At four carriers one SQP
iteration per tick flies as well as five in the twin's closed loop, but its plans are not converged on the rig's
flights, and at eight carriers only five iterations pass every bar (Section~\ref{sec:res_scale}). The eight-carrier
chain ran slower than real time, so it shows that the plan and the physics compose, not that the rig's laptop can fly
it.

\paragraph{The rim geometry}
The comparison with Quan et al.~\cite{quan2025swarm} turned on the attachment geometry rather than on the method's
code. With every cable at the ring's centre, the network held the ring's yaw to 0.1$^\circ$; with the cables on the
rim, as any payload with discrete attachment plates has them, the cables torque the ring in yaw and the network,
which has no yaw term, lets it swing (Section~\ref{sec:res_compare}). A planner that models each cable's attachment
point, as ours does, holds the yaw because the torque is in its model. For reconfiguration this matters doubly: a
release or a rejoin changes which plates carry the ring and so changes that torque.

\paragraph{Next on the rig}
The next rig flights test the attach of a newcomer on 1\,m cables, three carriers to four with a basket, with the
attach confirmed by motion capture before the fold-in; the four-carrier hover and paths on 1\,m cables; and the
planner's tick in flight with the speed-ups, which sets the real-time limit. The rejoin of a departed carrier, which
in the twin depends on the joint's weld, waits for a rig attach that has been shown repeatable.

\section{Conclusion}
\label{sec:conclusion}

We presented a planner that lets a team of quadrotors carry a cable-suspended ring through changes of its own
membership in flight. On the rig, an unannounced release of one of four carriers was detected in 0.42\,s, the team
was resized 0.09\,s later, and the ring tilted at most 6.5$^\circ$ before returning to its band in 1.8\,s. In the twin,
a departed carrier flew back to its own plate, welded within 1.5\,cm and was folded in with the ring tilting
1.3--3.4$^\circ$, and eight carriers took a larger ring down to four and back with a tilt peak of 3.35$^\circ$. Two SQP
iterations per tick converge for every fleet of up to four. Against a port of Quan et al.'s network method, our planner
held the ring's yaw where the network, lacking a yaw term for rim-attached cables, did not. The rig attach of a
newcomer on 1\,m cables is the next step \placeholder{result from 29 Oct}.
